\chapter*{Linear Models: Definitions, Properties, and Training}
\addcontentsline{toc}{chapter}{Linear Models: Definitions, Properties, and Training}

\section*{Introduction to Linear Models}

Linear models operate under the hypothesis that the target output can be well-described by a weighted linear combination of the input features. Despite their simplicity, they are fundamental, interpretable, and serve as the building blocks for more complex algorithms like neural networks.

\begin{definition}[Linear Model]
For an input vector $\xx = [x_1, x_2, \dots, x_d]^T$ with $d$ features, the linear model outputs a prediction $f_{\mathbf{w}}(\xx)$ computed as:
$$f_{\mathbf{w}}(\xx) = w_1x_1 + w_2x_2 + \dots + w_dx_d + w_0 = \mathbf{w}^T \xx + w_0$$
where $\mathbf{w} = [w_1, \dots, w_d]^T$ is the weight vector and $w_0$ is the bias (or intercept) term.
\end{definition}

\begin{remark}[Extended Feature Vector]
To simplify notation and avoid treating the bias $w_0$ separately, we typically augment the input vector with a constant $1$.
Let $\xx' = [1, x_1, \dots, x_d]^T$ and $\mathbf{w}' = [w_0, w_1, \dots, w_d]^T$. The prediction becomes a simple dot product:
$$f_{\mathbf{w}}(\xx) = (\mathbf{w}')^T \xx'$$
\end{remark}

\section*{Classification vs. Regression}

The interpretation of the linear output $f_{\mathbf{w}}(\xx)$ depends on the task:

\begin{description}
    \item[Regression:] The target output is a continuous curve fitting the data. The prediction is simply the raw output of the linear model: $\hat{y} = f_{\mathbf{w}}(\xx)$.
    \item[Classification:] The target output is a linear decision boundary (a hyperplane) that separates classes. For binary classification ($y \in \{-1, +1\}$), the prediction is the sign of the model's output: $\hat{y} = \text{sgn}(f_{\mathbf{w}}(\xx))$.
\end{description}

\begin{definition}[Functional Margin]
In binary classification, the \textbf{functional margin} is defined as $y \cdot f_{\mathbf{w}}(\xx)$, where $y \in \{-1, +1\}$ is the true label.
\begin{itemize}
    \item If $y \cdot f_{\mathbf{w}}(\xx) > 0$, the prediction is \textbf{correct} (the sign of the prediction matches the true label).
    \item If $y \cdot f_{\mathbf{w}}(\xx) < 0$, the prediction is \textbf{incorrect}.
\end{itemize}
The magnitude of the margin $|f_{\mathbf{w}}(\xx)|$ represents the \textit{confidence} of the prediction.
\end{definition}

\section*{Non-Linear Feature Transformations}

What if the data is not linearly separable or linearly distributed? We can still use linear models by applying a \textbf{feature transformation}.

\begin{remark}
We can map the original feature vector $\xx$ into a higher-dimensional space using a non-linear function $\phi(\xx)$. For example, if $\xx = [x_1, x_2]^T$, we could transform it into polynomial features:
$$\phi(\xx) = [1, x_1, x_2, x_1^2, x_2^2, x_1x_2]^T$$
The model $f_{\mathbf{w}}(\xx) = \mathbf{w}^T \phi(\xx)$ is \textbf{linear in the transformed parameter space}, but represents a \textbf{non-linear boundary/curve in the original input space}.
\end{remark}

\section*{Training: Empirical Risk Minimization}

Training a linear model means finding the optimal weights $\mathbf{w}^*$ that minimize the error over the training dataset. This is formalized as Empirical Risk Minimization (ERM).

\begin{definition}[Empirical Risk Minimization]
The goal is to find $\mathbf{w}$ that minimizes the average loss over the $N$ training examples:
$$\mathbf{w}^* = \arg\min_{\mathbf{w}} \frac{1}{N} \sum_{i=1}^N \ell\left(f_{\mathbf{w}}(\xx^{(i)}), y^{(i)}\right)$$
where $\ell(\hat{y}, y)$ is the \textbf{loss function} measuring the penalty for a bad prediction.
\end{definition}

\subsection*{Common Loss Functions}

\textbf{For Regression ($y \in \mathbb{R}$):}
\begin{description}
    \item[Squared Loss ($L_2$ Loss):] $\ell(\hat{y}, y) = (y - \hat{y})^2$. Strongly penalizes large outliers. Used in Ordinary Least Squares (OLS).
    \item[Absolute Loss ($L_1$ Loss):] $\ell(\hat{y}, y) = |y - \hat{y}|$. More robust to outliers.
\end{description}

\textbf{For Classification ($y \in \{-1, +1\}$, margin $\mu = y f_{\mathbf{w}}(\xx)$):}
\begin{description}
    \item[0-1 Loss:] $\ell(\mu) = 1$ if $\mu < 0$, else $0$. Ideal but computationally intractable (non-differentiable step function).
    \item[Hinge Loss:] $\ell(\mu) = \max(0, 1 - \mu)$. Used in Support Vector Machines (SVM).
    \item[Logistic Loss:] $\ell(\mu) = \log(1 + e^{-\mu})$. Smooth, differentiable convex surrogate. Used in Logistic Regression.
    \item[Exponential Loss:] $\ell(\mu) = e^{-\mu}$. Highly sensitive to misclassifications. Used in AdaBoost.
\end{description}

\section*{Regularization}

To prevent overfitting, especially when working with high-dimensional feature spaces (like expanded polynomials), we add a \textbf{regularization penalty} $R(\mathbf{w})$ to the loss function. The new objective becomes:
$$\min_{\mathbf{w}} \left( L_{\text{emp}}(\mathbf{w}) + \lambda R(\mathbf{w}) \right)$$
where $\lambda > 0$ controls the strength of the penalty.

\begin{description}
    \item[$\mathbf{L_2}$ Regularization (Ridge):] $R(\mathbf{w}) = ||\mathbf{w}||_2^2 = \sum w_j^2$. Penalizes large weights, driving them closer to zero, which shrinks variance and creates smoother decision functions.
    \item[$\mathbf{L_1}$ Regularization (Lasso):] $R(\mathbf{w}) = ||\mathbf{w}||_1 = \sum |w_j|$. Induces \textbf{sparsity}, meaning it forces less important feature weights exactly to zero, performing intrinsic feature selection.
\end{description}

\section*{Optimization: Finding the Optimal Weights}

How do we actually solve the $\arg\min$ problem?

\begin{enumerate}
    \item \textbf{Analytical Solution (Closed-form):} Set the gradient to zero ($\nabla_{\mathbf{w}} L(\mathbf{w}) = \mathbf{0}$) and solve for $\mathbf{w}$. This is possible for some models like unregularized or $L_2$-regularized linear regression.
    \item \textbf{Numerical Optimization (Gradient Descent):} For loss functions without closed-form solutions (like Logistic Loss), we use iterative methods. We compute the gradient $\nabla_{\mathbf{w}} L(\mathbf{w})$ and take steps in the opposite direction.
\end{enumerate}

\begin{definition}[Gradient Descent Update Rule]
At each iteration $t$, the weights are updated using the rule:
$$\mathbf{w}^{(t+1)} = \mathbf{w}^{(t)} - \eta \nabla_{\mathbf{w}} L(\mathbf{w}^{(t)})$$
where $\eta$ is the \textbf{learning rate} (step size), determining how far we move along the slope of the error surface at each step.
\end{definition}
